\lecture{7}
\title{Kernels, Evaluations, and Hyperparameter Learning}

\begin{document}

\subsection{Kernels: Properties and Examples}

\begin{problem}
    Suppose our input space is $\mathcal{X} = \mathbb{R}^2$,
    and our kernel is $k: \mathbb{R}^2 \times \mathbb{R}^2 \to \mathbb{R}$
    via $k(a, b) := (a^{\top}b)^2$.
    What feature map $\phi$ is represented by this kernel,
    and what is the implied feature dimension?
\end{problem}

\begin{solution}
    Given inputs $a = \begin{bmatrix} a_1 & a_2 \end{bmatrix}^{\top}$
    and $b = \begin{bmatrix} b_1 & b_2 \end{bmatrix}^{\top}$,
    the kernel can be expressed as:
    \[ k(a, b) = a_1^2b_1^2 + 2a_1a_2b_1b_2 + a_2^2b_2^2
    = \begin{bmatrix}
    a_1^2 \\ \sqrt{2} a_1a_2 \\ a_2^2
    \end{bmatrix}^{\top}
    \begin{bmatrix}
    b_1^2 \\ \sqrt{2} b_1b_2 \\ b_2^2
    \end{bmatrix}. \]
    Having ``decoupled'' the kernel in this way,
    it is now clear that the feature map is
    $\phi(x) := \begin{bmatrix}
        x_1^2 & \sqrt{2}x_1x_2 & x_2^2
    \end{bmatrix}^{\top}$,
    and the implied feature dimension is $3$.
    (More generally, the implied feature dimension
    is the ``number of terms'' in the kernel.)
\end{solution}

To really emphasize how much cleaner kernels are to work with
than features, here's a much larger example.

\begin{problem}
    Suppose our input space is $\mathcal{X} = \mathbb{R}^D$,
    and our kernel is $k: \mathbb{R}^D \times \mathbb{R}^D \to \mathbb{R}$
    via $k(a, b) := (a^{\top}b)^M$.
    What is the implied feature dimension?
    Compare the difficulty of computing the matrices $K$ versus $\Phi$
    (using the notation from \href{/6.790/lectures/6}{Lecture 6}).
\end{problem}

\begin{solution}
    The feature dimension is the ``number of terms'' in
    the expansion of $k(a, b) := (a^{\top}b)^{M}$.
    This is equivalent to the number of monomials of degree $M$
    in $D$ variables, which (by Stars and Bars) is:
    \[ \text{feature dimension} = \binom{M + D - 1}{D - 1}. \]
    Now for the computational difficulty comparisons.
    \begin{itemize}
        \item Computing $K$ costs $O(D \cdot N^2)$.
        \begin{itemize}
            \item We need the $N^2$ values $k(x^{(i)}, x^{(j)})$
        across all $N^2$ pairs of inputs $\{x^{(n)}\}_{n = 1}^N$,
        and computing each value takes $O(D)$.
        \end{itemize}
        \item Computing $\Phi$ costs $\Omega\left(
            N \cdot \binom{M + D - 1}{D - 1}
        \right)$.
        \begin{itemize}
            \item This bound comes from the space of $\Phi$, alone:
            the dimensions of $\Phi$ are
            $N \times \binom{M + D - 1}{D - 1}$.
        \end{itemize}
    \end{itemize}
    Evidently, it is much more convenient to work with kernels
    than it is to work with features.
\end{solution}

And it turns out that kernels do everything that features can do.

\begin{theorem}{Mercer's Theorem}
    If, for every choice of inputs $\{x^{(n)}\}_{n = 1}^N$,
    the kernel matrix $K_{i, j} := k(x^{(i)}, x^{(j)})$
    is symmetric and positive semidefinite, and
    the kernel $k\colon \mathcal{X} \times \mathcal{X} \to \mathbb{R}$
    is continuous,
    then $k$ corresponds to some feature map $\phi$.
\end{theorem}

The proof is fairly nontrivial, so we omit it.
The upshot is that we can now declare any kernel we'd like
(so long as it satisfies the conditions of Mercer's theorem),
and we need not have any concern for what features are involved.

\begin{remark}
    If we can choose kernels arbitrarily,
    that raises the question:
    what makes one choice of kernel better than another?

    Here's one possible answer,
    from the ``feature'' perspective.

    \begin{quote}
        \textbf{Answer 1.}
        The kernel ultimately describes a collection of $P$ features,
        and a ridge regression models the relationship between
        $\mathcal{X}$ and $\mathcal{Y}$ via
        a linear combination of those features,
        with preference towards linear combinations
        with smaller coefficients.

        Therefore, a good kernel should imply a collection of features
        that is capable of modeling the relationship
        between $\mathcal{X}$ and $\mathcal{Y}$
        using relatively small coefficients.
    \end{quote}

    Here's another possible answer,
    from the ``hypothesis'' perspective.

    \begin{quote}
        \textbf{Answer 2.}
        Recall from
        \href{/6.790/lectures/6/#kernels-unlimited-features}{Lecture 6}
        that, ultimately, our hypothesis $h\colon \mathcal{X} \to \mathbb{R}$
        looks like
        $h(x) = Y^{\top}(\lambda I_N + K)^{-1}k(x)$.
        Expressed component-wise, that looks like:
        \[ h(x)
        = \left[(\lambda I_N + K)^{-1}Y\right]^{\top} k(x)
        = \sum_{n = 1}^N \alpha_n \cdot k(x^{(n)}, x)
        ~~~ \text{ where } ~~~
        \alpha = (\lambda I_N + K)^{-1}Y \in \mathbb{R}^N.  \]
        In other words, $h(x)$ is a weighted sum of the $\alpha_n$,
        with weights given by $k(x^{(n)}, x)$.
        This brings out a sense in which $k(a, b)$
        reflects a \emph{similarity score} between $a$ and $b$.
        The more similar $x^{(n)}$ is to $x$,
        the more impact $\alpha_n$ has on $h(x)$.

        The kernel can also be seen as a similarity score
        by recalling that $k(a, b) := \phi^{\top}(a) \phi(b)$
        is just the dot product of $\phi(a)$ and $\phi(b)$
        in the feature space $\mathbb{R}^P$.
        Indeed, if we normalize the kernel
        so that $k(x, x) \equiv 1$, then we have
        \[ \|\phi(a) - \phi(b)\|^2 = 2 - 2 k(a, b) \]
        by the Law of Cosines.
        The similarity score $k(a, b)$
        captures the similarity of $\phi(a)$ and $\phi(b)$
        as positions in the feature space.
    \end{quote}
\end{remark}

\textbf{Definition.}
One common choice of kernel is the
\textbf{Gaussian} or \textbf{RBF} kernel, given by:
\[ k(a, b) := \exp\left(
    -\frac{\|a - b\|^2}{2\ell^2}
\right) ~ \text{ for some } \ell > 0. \]
Note that this kernel has an
\emph{infinite-dimensional feature map}.
(All of the work done in the previous lecture
works just as well in infinite-dimensional vector spaces,
so there's nothing to be worried about.)

\subsection{Recap: Two Approaches to Regression}

Let's summarize all that we've discussed about regression so far.

Given some data $\{(x^{(n)}, y^{(n)})\}_{n = 1}^N$,
there are two approaches to learning a hypothesis
$h \colon \mathcal{X} \to \mathcal{Y}$.
\[
    h(x) := \theta^{\top}\phi(x)
    ~ \text{ where } ~
    \theta = \Phi^{\top}(\lambda I_N + K)^{-1}Y
    ~~~~~ \parallel ~~~~~
    h(x) := Y^{\top}(\lambda I_N + K)^{-1}k(x).
\]
These two approaches conceptually differ in the following way:
\begin{itemize}
    \item The former involves learning a parameter $\theta$
    in the feature space described by $\phi(x)$.
    \item The latter is a nonparametric definition of
    a function in a sample space whose behavior
    depends on the kernel $k(x^{(i)}, x^{(j)})$.
\end{itemize}
Besides their conceptual differences, however,
these two approaches are the same:
the hypotheses $h(x)$ they generate are identical.

\subsection{Learning Algorithms and Hyperparameters}

There are two types of objects
that appear in the previous section,
and it's important to distinguish them clearly.
\begin{itemize}
    \item A \textbf{hypothesis} (or \textbf{decision rule})
    $h \colon \mathcal{X} \to \mathcal{Y}$
    (recalling \href{/6.790/lectures/1/#setting-up}{Lecture 1})
    takes an input $x$ and returns a hypothesis $h(x)$.
    \item A \textbf{learning algorithm} $\mathcal{A}$
    takes training data $\mathcal{D} = \{(x^{(n)}, y^{(n)})\}_{n = 1}^N$
    and returns a hypothesis $h = \mathcal{A}(\mathcal{D})$.
\end{itemize}
In the previous section,
the two approaches described are both
(the same) \emph{learning algorithm},
whereas the results $h(x)$ produced by these approaches
are \emph{hypotheses} (or \emph{decision rules}).

Notably, the behavior of the hypothesis $h(x)$,
as presented by the first of the two approaches,
depends on a \textbf{parameter} $\theta$.
In contrast, the behavior of the
learning algorithm $\mathcal{A}_{\lambda}$
depends on a \textbf{hyperparameter} $\lambda$.

\subsection{Evaluating Learning Algorithms}

Recall that in
\href{/6.790/lectures/1/#the-optimal-hypothesis-given-the-distribution}{Lecture 1}
we defined the \textbf{risk}
of a hypothesis $h$ to be:
\[ \mathrm{risk}(h) := \mathbb{E}\left[L(Y, h(X))\right], \]
where $L \colon \mathcal{Y} \times \mathcal{Y} \to \mathbb{R}$
is a loss function of our choice.
Note that risk is a property of a hypothesis,
not of a learning algorithm---but
the risk of a hypothesis
can be controlled by carefully setting the
hyperparameter of its learning algorithm.

Here's the focus of this section:

\begin{quote}
    \textbf{Question.} How do we evaluate
    a learning algorithm $\mathcal{A}_{\lambda}$
    to determine the best values
    for its hyperparameters?
\end{quote}

As with bootstrapping (discussed in
\href{/18.650/lectures/13/#bootstrapping}{Lecture 13} of 18.650),
the idea is to empirically estimate
the expected risk of the hypothesis produced,
using our own training data as validation data.
Stated more precisely:
\begin{enumerate}
    \item Randomly partition the data $\mathcal{D}$
    into a training set $\mathcal{D}_T$
    and a validation set $\mathcal{D}_V$.
    \item Run the learning algorithm on the training set
    to obtain a hypothesis
    $h_{\mathcal{D}_T, \lambda} := \mathcal{A}_{\lambda}(\mathcal{D}_T)$.
    \item Estimate the risk of this hypothesis
    to be the average loss across the validation data $\mathcal{D}_V$.
    \[ \widehat{\mathrm{risk}}(h_{\mathcal{D}_T, \lambda})
    = \frac{1}{|\mathcal{D}_V|} \sum_{(x, y) \in \mathcal{D}_V}
    L(y, h_{\mathcal{D}_T, \lambda}(x)). \]
    \item Repeat the previous three steps across many trials,
    and average the results from step 3 across all trials.
\end{enumerate}
The Law of Large Numbers says that
the risk estimates in step 3
will approach the true risk for large $|\mathcal{D}_V|$,
assuming all data is sampled i.i.d. from the same distribution.
You can then quantify the (epistemic) uncertainty
in the estimate of expected risk produced by step 4 by using
basic statistics as well.

\begin{remark}
    Pay attention to the difference between
    \emph{risk} and \emph{expected risk}.
    \begin{itemize}
        \item An iteration of step 3 produces an estimate of
        the \emph{risk} of $h_{\mathcal{D}_T, \lambda}$.
        \item Step 4 estimates the \emph{expected risk}
        by averaging the risk estimates from step 3
        across multiple different hypotheses
        $\{h_{\mathcal{D}_T^{(b)}, \lambda}\}_{b = 1}^B$,
        each trained on its own training set $\mathcal{D}_T^{(b)}$
        and evaluated on its own validation set $\mathcal{D}_V^{(b)}$.
    \end{itemize}
\end{remark}

Whereas the parameter $\theta$ is
evaluated by its risk
and learned by the learning algorithm,
the hyperparameter $\lambda$ is
evaluated by its expected risk
and learned by us.

\end{document}